% ============================================
% chapters/before_departure.tex
\section{行前准备}
% ============================================

% --------------------
% Subsection: 申请流程总览（原 section 下调一级）
% --------------------
\subsection{申请流程总览}

\subsubsection{待办事项总览}

% 【子无编号小节：拿到学校 offer 后 → 下调为 paragraph*】
\paragraph*{拿到学校 offer 后}

% 灰条跨页不断裂设置（若灰条环境在下方，提前设置避免报错）
\tcbset{before skip=4pt, after skip=4pt}

% 待办事项清单（入学前）
\begin{itemize}[label=•, itemsep=2pt, leftmargin=2em]
  \item 申请住宿
  \item 缴纳学费 \& 保险
  \item 换 uncon（仅应届毕业生）
  \item 办理签证
  \item 购买机票
\end{itemize}

% 【子无编号小节：入境丹麦后 → 下调为 paragraph*】
\paragraph*{入境丹麦后}

% 待办事项清单（入境后）
\begin{itemize}[label=•, itemsep=2pt, leftmargin=2em]
  \item 按居住地规则申请 CPR，并办理 MitID、Digital Post 和黄卡；居留卡由 SIRI 或丹麦移民局制发
  \item 办理丹麦银行卡
\end{itemize}

% --------------------
% Subsection: 录取后 Checklist（原 section 下调一级）
% --------------------
\subsection{录取后 Checklist}

\subsubsection{申请住宿}
详见《留学生手册》“租房信息 \& 宿舍介绍”板块。

\subsubsection{缴纳学费 \& 保险}

% 灰条提示框（带标题）
\begin{grayleftbox}\textbf{学长学姐碎碎念}
\begin{itemize}
  \item “我在北京分行中午办的中行一类卡，需要预约，但当时很好约。”
  \item “拿到一类卡后可以买外汇，然后转到自己名下的万事达或 Visa 储蓄卡／信用卡。”
  \item “建议申请一张万事达储蓄卡（通常比 Visa 更易申请），用来境外消费和交学费。”
  \item “北京分行工作人员查看 offer 后给我开的额度很大，非常方便！”
  \item “该储蓄卡开通 POS 机免密并设置额度后几乎可当信用卡用，但每笔消费会多锁 2\% 金额，且绑卡不太方便。”
  \item “外币信用卡（如中行卓隽）可参考小红书教程申请，一张基本就够用！”
\end{itemize}
\end{grayleftbox}

\subsubsection{办理签证}
详见《留学生手册》“签证与居留许可”板块。

%\subsection{机票 \& 行李／转运}

\subsection{常见航空公司及飞行体验}

% 使用表格展示航司信息（排版紧凑，灰白条）
\begin{table}[H]
  \centering
  \scriptsize
  \renewcommand{\arraystretch}{1.2}
  \setlength{\tabcolsep}{4pt}
  \rowcolors{2}{gray!10}{white}
  \begin{tabularx}{\textwidth}{
      p{0.10\textwidth}  % 航空公司
      X  % 航线特点
      p{0.22\textwidth}  % 行李政策
      p{0.22\textwidth}  % 其他备注
    }
    \toprule
    \textbf{航空公司} & \textbf{航线特点} & \textbf{行李政策} & \textbf{其他备注} \\
    \midrule
    中国国航 &
      北京首都机场直飞哥本哈根；寒暑假票价上涨，8 月末单程约 4–5 千元；机上有 Wi-Fi &
      留学生票可托运 2 × 23 kg &
      -- \\
    泰国航空 &
      CPH出发飞 10.5 h 到曼谷；适合家在西南的同学 &
      无留学生票，只托运 1 × 23 kg &
      无 Wi-Fi；机餐好；曼谷需重新安检，勿在 CPH机场买酒 \\
    新加坡航空 &
      新加坡—哥本哈根约 13 h &
      参考官网 &
      -- \\
    汉莎航空 &
      非直达航班，整体飞行时间较长，经济舱空间一般 &
      留学生票（博士也适用）可享更多免费行李额 &
      -- \\
    \bottomrule
  \end{tabularx}
\end{table}

% --------------------
% Subsection: 入境材料（原 section 下调一级）
% --------------------
\subsection{入境材料}

% 入境材料清单
\begin{itemize}[label=•, itemsep=2pt, leftmargin=2em]
  \item 护照（含签证页）
  \item \textbf{Residence permit letter}：居留许可获批信；在居留卡尚未收到时是重要的许可证明，请保存电子版并按出入境要求携带所需材料
  \item 学校 Offer
  \item 租房合同
\end{itemize}

% 灰条提示框（带标题）
\begin{grayleftbox}\textbf{提示: }
建议所有材料携带原件及复印件各 1 份。
\end{grayleftbox}

% ----------------------- End of file -----------------------
